% Source: https://artofproblemsolving.com/community/c3h1695868_2018_nats_mc_questions_sprint_13_14_23
% Stable identity: aops-topic-1695868-q14
Two faces sum to $12$ only as $6+6$, so the event is ``at least two of the four dice show $6$''. Use the complement: no $6$ has probability $\left(\dfrac{5}{6}\right)^4=\dfrac{625}{1296}$, and exactly one $6$ has probability $4\cdot\dfrac{1}{6}\cdot\left(\dfrac{5}{6}\right)^3=\dfrac{500}{1296}$. Hence
\[
P=1-\dfrac{625+500}{1296}=\dfrac{171}{1296}=\boxed{\dfrac{19}{144}}.
\]
(The poster's $\dfrac{211}{1296}$ is too large by $\dfrac{40}{1296}$; the count $171=\dbinom42\cdot25+\dbinom43\cdot5+1$ of outcomes with two, three or four sixes confirms the answer.)
